\section {Frogs over the x axis (Combinatorics)}

Frogs sit on points of the plane with whole coordinates. A frog may jump over a frog that is right next to it
(north, south, east or west) and land on the empty point right behind it, at a distance of $2$ from where it started.
The frog that was jumped over is removed.

At the start there is a frog on every point $(x,y)$ with $y \le 0$, and all the points with $y > 0$ are empty.
How far north can a frog get after a finite number of jumps? Specifically, prove that no frog can reach the line $y=5$.

\subsection{Solution}

Assume a frog can reach some point on the line $y=5$; call this point the target.
Let $\sigma = \frac{\sqrt{5}-1}{2} \approx 0.618$. It is the positive root of $\sigma^2 + \sigma = 1$.
Give a frog that sits at distance $d$ from the target the weight $\sigma^d$, where the distance is measured along the grid,
$|x-x_0|+|y-y_0|$, and let the weight of a set of frogs be the sum of their weights.
A frog on the target has weight $1$.

Every jump replaces two frogs by one. There are only three kinds of jumps:
\begin{enumerate}
\item A frog at distance $d+2$ jumps over a frog at distance $d+1$ and lands at distance $d$.
The weight changes by $\sigma^d - \sigma^{d+1} - \sigma^{d+2} = \sigma^d(1-\sigma-\sigma^2) = 0$.
\item A frog at distance $d$ jumps over a frog at distance $d+1$ to distance $d+2$.
The weight changes by $\sigma^{d+2} - \sigma^d - \sigma^{d+1} < 0$.
\item A frog lands at the same distance it started from (for example when jumping across the column of the target).
The weight changes by minus the weight of the frog that was jumped over, which is negative.
\end{enumerate}
So the weight of the frogs never grows.

Now let's compute the weight of all the frogs at the start. Put the target at $(0,5)$.
The frogs on the line $y=0$ are at distances $5, 6, 6, 7, 7, \ldots$ from the target, so their weight is
\begin{center}
$\sigma^5 + 2\sigma^6 + 2\sigma^7 + \ldots = \sigma^5 + \frac{2\sigma^6}{1-\sigma} = \sigma^5 + 2\sigma^4$
\end{center}
using $1-\sigma = \sigma^2$. Every line further south is one step further away, so the weight of all the frogs is
\begin{center}
$\frac{\sigma^5 + 2\sigma^4}{1-\sigma} = \sigma^3 + 2\sigma^2 = (2\sigma - 1) + 2(1-\sigma) = 1$
\end{center}
using $\sigma^2 = 1-\sigma$ and $\sigma^3 = \sigma - \sigma^2 = 2\sigma - 1$.

A finite number of jumps involves only a finite number of frogs. Their weight at the start is a finite part of a sum of
positive numbers that adds up to $1$, so it is strictly less than $1$. The jumps never increase it, so at the end it is still less
than $1$, and no frog can be on the target.

On the other hand, frogs can reach the line $y=4$ (it takes 20 frogs), so $4$ is as far north as they can get.
This riddle is known as Conway's soldiers.
